% GENERATED FILE. Edit canonical lessons, then run npm run generate:books.
% canonical-source-hash: fnv1a64:eef2144ef076b5ec
% canonical-lessons: HI-C331-vrat, HI-C331-arti, HI-C331-bhajan, HI-C331-prasad, HI-C331-pakvan

\chapter{Fasting and Offerings}
\label{ch:hi-fasting-and-offerings}
\hlchaptermodality{331}

\begin{chapteropening}
\textbf{By the end of this chapter:} \emph{I can say how a family kept a festival day.}

Complete the last lesson of chapter 331: I can say how a family kept a festival day.
\end{chapteropening}

\section[\texorpdfstring{\hi{व्रत} (vrat)}{व्रत (vrat)}]{\hi{व्रत} (vrat) — a fast; a vow}

\label{lesson:HI-C331-vrat}



\begin{quote}
Before the new one: say the Hindi for a festive procession, then the Hindi for a devotee.
\end{quote}

\subsection*{You'll want to know: \hi{व्रत}}
\textbf{\hi{व्रत}} — \emph{vrat} — \textquotedblleft{}a fast; a vow\textquotedblright{}.

\begin{quote}
\hi{दादी} \hi{ने} \hi{पूरे} \hi{दिन} \hi{व्रत} \hi{रखा।}
\end{quote}

\begin{quote}
\emph{dādī ne pūre din vrat rakhā.}
\end{quote}

\begin{quote}
Grandmother fasted all day.
\end{quote}

Masculine.

\begin{cousinweb}
\textbf{\hi{व्रत}} (\emph{vrat}) is Sanskrit \emph{vrata}, \textquotedblleft{}a vow, a rule of conduct\textquotedblright{} (Monier-Williams). \textbf{\hi{व्रत} \hi{रखना}} is to keep a fast, a vow to go without food.
\end{cousinweb}

\begin{grammarlens}[title={What you've built}]
One more word for telling what happened.
\end{grammarlens}

\subsection*{Your turn}
\begin{itemize}
\raggedright
  \item \emph{Say it:} \emph{vrat}
  \item \emph{Say it:} \emph{vrat}, once more
  \item \emph{Say it:} \emph{dādī ne pūre din vrat rakhā.} — Grandmother fasted all day.
  \item \emph{Recall:} say the Hindi for a festive procession, then the Hindi for a devotee, then say \emph{vrat} again
\end{itemize}

\begin{tcolorbox}[breakable,colback=teal!4,colframe=teal!35!black,title={Before you move on}]
What does \hi{व्रत} mean? (\textquotedblleft{}A fast; a vow\textquotedblright{}.) Say it once more.
\end{tcolorbox}
\section[\texorpdfstring{\hi{आरती} (ārtī)}{आरती (ārtī)}]{\hi{आरती} (ārtī) — an arati, worship with a lamp}

\label{lesson:HI-C331-arti}



\begin{quote}
Before the new one: say the Hindi for a devotee, then the Hindi for a fast.
\end{quote}

\subsection*{You'll want to know: \hi{आरती}}
\textbf{\hi{आरती}} — \emph{ārtī} — \textquotedblleft{}an arati, worship with a lamp\textquotedblright{}.

\begin{quote}
\hi{शाम} \hi{को} \hi{मंदिर} \hi{में} \hi{आरती} \hi{हुई।}
\end{quote}

\begin{quote}
\emph{śām ko mandir meṁ ārtī huī.}
\end{quote}

\begin{quote}
In the evening there was an arati at the temple.
\end{quote}

Feminine.

\begin{cousinweb}
\textbf{\hi{आरती}} (\emph{ārtī}) goes back to Sanskrit \emph{ārātrika} (CDIAL 1315): a lamp waved in circles before a god while a song is sung.
\end{cousinweb}

\begin{grammarlens}[title={What you've built}]
One more word for telling what happened.
\end{grammarlens}

\subsection*{Your turn}
\begin{itemize}
\raggedright
  \item \emph{Say it:} \emph{ārtī}
  \item \emph{Say it:} \emph{ārtī}, once more
  \item \emph{Say it:} \emph{śām ko mandir meṁ ārtī huī.} — In the evening there was an arati at the temple.
  \item \emph{Recall:} say the Hindi for a devotee, then the Hindi for a fast, then say \emph{ārtī} again
\end{itemize}

\begin{tcolorbox}[breakable,colback=teal!4,colframe=teal!35!black,title={Before you move on}]
What does \hi{आरती} mean? (\textquotedblleft{}An arati, worship with a lamp\textquotedblright{}.) Say it once more.
\end{tcolorbox}
\section[\texorpdfstring{\hi{भजन} (bhajan)}{भजन (bhajan)}]{\hi{भजन} (bhajan) — a hymn; a devotional song}

\label{lesson:HI-C331-bhajan}



\begin{quote}
Before the new one: say the Hindi for a fast, then the Hindi for an arati, worship with a lamp.
\end{quote}

\subsection*{You'll want to know: \hi{भजन}}
\textbf{\hi{भजन}} — \emph{bhajan} — \textquotedblleft{}a hymn; a devotional song\textquotedblright{}.

\begin{quote}
\hi{आरती} \hi{के} \hi{बाद} \hi{सबने} \hi{भजन} \hi{गाए।}
\end{quote}

\begin{quote}
\emph{ārtī ke bād sabne bhajan gāe.}
\end{quote}

\begin{quote}
After the arati everyone sang hymns.
\end{quote}

Masculine.

\begin{cousinweb}
\textbf{\hi{भजन}} (\emph{bhajan}) is Sanskrit \emph{bhajana}, \textquotedblleft{}worship, adoration\textquotedblright{} (Monier-Williams), the same family as \textbf{\hi{भक्त}}: a song of worship.
\end{cousinweb}

\begin{grammarlens}[title={What you've built}]
One more word for telling what happened.
\end{grammarlens}

\subsection*{Your turn}
\begin{itemize}
\raggedright
  \item \emph{Say it:} \emph{bhajan}
  \item \emph{Say it:} \emph{bhajan}, once more
  \item \emph{Say it:} \emph{ārtī ke bād sabne bhajan gāe.} — After the arati everyone sang hymns.
  \item \emph{Recall:} say the Hindi for a fast, then the Hindi for an arati, worship with a lamp, then say \emph{bhajan} again
\end{itemize}

\begin{tcolorbox}[breakable,colback=teal!4,colframe=teal!35!black,title={Before you move on}]
What does \hi{भजन} mean? (\textquotedblleft{}A hymn; a devotional song\textquotedblright{}.) Say it once more.
\end{tcolorbox}
\section[\texorpdfstring{\hi{प्रसाद} (prasād)}{प्रसाद (prasād)}]{\hi{प्रसाद} (prasād) — a blessed offering (of food)}

\label{lesson:HI-C331-prasad}



\begin{quote}
Before the new one: say the Hindi for an arati, worship with a lamp, then the Hindi for a hymn.
\end{quote}

\subsection*{You'll want to know: \hi{प्रसाद}}
\textbf{\hi{प्रसाद}} — \emph{prasād} — \textquotedblleft{}a blessed offering (of food)\textquotedblright{}.

\begin{quote}
\hi{पुजारी} \hi{ने} \hi{सबको} \hi{प्रसाद} \hi{बाँटा।}
\end{quote}

\begin{quote}
\emph{pujārī ne sabko prasād bā̃ṭā.}
\end{quote}

\begin{quote}
The priest handed out the blessed offering to everyone.
\end{quote}

Masculine.

\begin{cousinweb}
\textbf{\hi{प्रसाद}} (\emph{prasād}) is Sanskrit \emph{prasāda}, \textquotedblleft{}grace, favour\textquotedblright{} (Monier-Williams): food first offered to a god and then shared out as a blessing.
\end{cousinweb}

\begin{grammarlens}[title={What you've built}]
One more word for telling what happened.
\end{grammarlens}

\subsection*{Your turn}
\begin{itemize}
\raggedright
  \item \emph{Say it:} \emph{prasād}
  \item \emph{Say it:} \emph{prasād}, once more
  \item \emph{Say it:} \emph{pujārī ne sabko prasād bā̃ṭā.} — The priest handed out the blessed offering to everyone.
  \item \emph{Recall:} say the Hindi for an arati, worship with a lamp, then the Hindi for a hymn, then say \emph{prasād} again
\end{itemize}

\begin{tcolorbox}[breakable,colback=teal!4,colframe=teal!35!black,title={Before you move on}]
What does \hi{प्रसाद} mean? (\textquotedblleft{}A blessed offering (of food)\textquotedblright{}.) Say it once more.
\end{tcolorbox}
\section[\texorpdfstring{\hi{पकवान} (pakvān)}{पकवान (pakvān)}]{\hi{पकवान} (pakvān) — a festive dish; a delicacy}

\label{lesson:HI-C331-pakvan}



\begin{quote}
Before the new one: say the Hindi for a hymn, then the Hindi for a blessed offering (of food).
\end{quote}

\subsection*{You'll want to know: \hi{पकवान}}
\textbf{\hi{पकवान}} — \emph{pakvān} — \textquotedblleft{}a festive dish; a delicacy\textquotedblright{}.

\begin{quote}
\hi{माँ} \hi{ने} \hi{त्योहार} \hi{के} \hi{लिए} \hi{कई} \hi{पकवान} \hi{बनाए।}
\end{quote}

\begin{quote}
\emph{māṁ ne tyohār ke liye kaī pakvān banāe.}
\end{quote}

\begin{quote}
Mother made several festive dishes for the festival.
\end{quote}

Masculine.

\begin{cousinweb}
\textbf{\hi{पकवान}} (\emph{pakvān}) is Sanskrit \emph{pakvānna}, \emph{pakva}, cooked, and \emph{anna}, food (Monier-Williams); in Hindi it is the rich, fried food of a feast day.
\end{cousinweb}

\begin{grammarlens}[title={What you've built}]
One more word for telling what happened.
\end{grammarlens}

\subsection*{Your turn}
\begin{itemize}
\raggedright
  \item \emph{Say it:} \emph{pakvān}
  \item \emph{Say it:} \emph{pakvān}, once more
  \item \emph{Say it:} \emph{māṁ ne tyohār ke liye kaī pakvān banāe.} — Mother made several festive dishes for the festival.
  \item \emph{Recall:} say the Hindi for a hymn, then the Hindi for a blessed offering (of food), then say \emph{pakvān} again
\end{itemize}

\begin{tcolorbox}[breakable,colback=teal!4,colframe=teal!35!black,title={Before you move on}]
What does \hi{पकवान} mean? (\textquotedblleft{}A festive dish; a delicacy\textquotedblright{}.) Say it once more.
\end{tcolorbox}
